\chapter*{General Introduction}
\addcontentsline{toc}{chapter}{General Introduction}

\section{Context of the Internship and Project}

This final-year engineering project was carried out at Africom Services, a Tunisian company founded in January 2005 by experienced telecommunications executives. The company operates in telecommunications engineering and network maintenance, with activities covering transmission systems, radio-electric networks, fiber-optic deployment, electronic security, and integrated energy/IT solutions.

Africom Services supports telecom operators, private companies, and public institutions in modernizing communication infrastructure. Its technical positioning is reinforced by a collaboration with Telecom Canada, which provides additional expertise in private networks, wireless communication, and digital infrastructure.

Within this context, the company identified a recurring operational challenge in field intervention management: intervention data were often captured through fragmented and partially manual workflows, making monitoring, validation, and reporting slower and less reliable than required for modern service delivery.

\section{Problem Statement and Motivation}

The management of IT equipment installation and maintenance interventions involves multiple stakeholders, distributed teams, and strict quality constraints. In practice, the existing process generated five major limitations:

\begin{enumerate}
\item Delays between intervention execution and report delivery.
\item Heterogeneous data capture quality across technicians and projects.
\item Limited real-time visibility for supervisors and coordinators.
\item Difficulty in enforcing standardized workflows and validation rules.
\item Administrative overhead in consolidating intervention outputs for clients.
\end{enumerate}

These limitations motivated the design of an integrated digital platform able to structure the complete intervention lifecycle from planning to execution, validation, and reporting.

\section{Project Objective}

The objective of this project is to design and implement \textbf{RM (Report Manager)}, a web and mobile solution dedicated to intervention management for IT equipment installation and maintenance activities.

The project pursues the following specific objectives:

\begin{itemize}
\item Digitize and standardize intervention workflows.
\item Enable role-based collaboration between technicians, supervisors, coordinators, and administrators.
\item Improve traceability through status control, audit information, and structured evidence management.
\item Automate the production of operational reports for decision-making and client communication.
\item Provide a scalable technical foundation aligned with industrial software engineering practices.
\end{itemize}

\section{Methodological Approach}

The project follows an agile SCRUM process based on incremental delivery. Development was organized into iterative sprints, each targeting a coherent functional scope and ending with validation activities.

This methodology was selected to:

\begin{itemize}
\item maintain continuous alignment with operational needs,
\item reduce technical and organizational risk through progressive delivery,
\item and integrate feedback loops for quality improvement.
\end{itemize}

In practical execution, the SCRUM framework was adapted to a single-developer delivery model with explicit control of backlog prioritization, sprint planning, and sprint review checkpoints. This adaptation preserved methodological rigor while remaining compatible with academic time constraints.

From a technical perspective, the implementation combines a backend service layer, a web interface, and a mobile client, supported by a secure and structured data management strategy.

\section{Report Organization}

The remainder of this report is organized as follows:

\begin{itemize}
\item \textbf{Chapter 1: General Presentation} introduces the host company, project environment, problematic, and objectives.
\item \textbf{Chapter 2: Project Planning} presents the SCRUM process, tools, and physical/logical architecture decisions.
\item \textbf{Chapter 3: Analysis and Requirements Specification} formalizes actors, role codes, functional and non-functional requirements, use cases, and backlog structure.
\item \textbf{Chapter 4: Sprint 1} details foundation and core infrastructure implementation.
\item \textbf{Chapter 5: Sprint 2} presents rapport request and dynamic form template realization.
\item \textbf{Chapter 6: Sprint 3} covers file management and geolocation tracking features.
\item \textbf{Chapter 7: Sprint 4} documents approval workflow, rejection logic, and status governance.
\item \textbf{Chapter 8: Sprint 5} presents reporting, export integration, and optimization mechanisms.
\item \textbf{General Conclusion} synthesizes achievements, limitations, and future work perspectives.
\end{itemize}

This structure ensures a coherent progression from contextual framing and requirement formalization to design rationale, implementation details, and critical evaluation of results.
